\chapter*{术语表}

\begin{description}
\item[算术强度] 操作数与数据移动字节数的比值，用来判断一个内核更可能受计算吞吐还是内存带宽限制。
\item[Roofline] 用计算峰值、内存带宽和算术强度估算性能上界的模型。
\item[SIMD] 一条指令处理多个数据 lane 的执行方式，是现代 CPU 单核数值吞吐的重要来源。
\item[分块] 把大问题切成适合缓存或寄存器的小块，以提高数据复用。
\item[Packing] 把矩阵块复制成连续、对齐、适合微内核访问的临时布局。
\item[微内核] GEMM 或卷积中最内层的架构相关计算核心，通常直接围绕寄存器和 SIMD 指令设计。
\item[KV Cache] Transformer decode 阶段保存历史 key 和 value 的缓存，用于避免重复计算历史 token。
\item[量化] 用低位宽整数或低精度浮点表示权重、激活或缓存，并通过 scale、zero point 等元数据近似原始实数。
\item[NUMA] 非统一内存访问结构，不同核心访问不同内存节点的成本不同。
\item[算子融合] 把多个相邻算子合并执行，减少中间张量读写和调度开销。
\end{description}
